\documentclass[11pt]{article}
\usepackage[a4paper,margin=1in]{geometry}
\usepackage{amsmath,amsthm,amssymb}
\usepackage{mathtools}

\newtheorem{theorem}{Theorem}[section]
\newtheorem{proposition}[theorem]{Proposition}
\newtheorem{lemma}[theorem]{Lemma}
\newtheorem{corollary}[theorem]{Corollary}
\theoremstyle{definition}
\newtheorem{remark}[theorem]{Remark}

\newcommand{\Om}{\Omega}
\newcommand{\rstar}{\rho_*}
\newcommand{\Hn}{\mathcal H^{n-1}}
\newcommand{\X}{\mathcal X}

\title{Self-Truncated Velocity Defect\\
       Replacing the Pointwise Lower Gradient Hypothesis}
\author{}
\date{}

\begin{document}
\maketitle

\begin{abstract}
The original Plan 2 metric trace used a pointwise lower bound
$|\nabla u|\ge m_0$ on every reduced-boundary level surface.  That
former global lower-gradient hypothesis is too strong: tiny components
can carry arbitrarily small gradients.  This note records the
replacement.  On a.e. level whose perimeter is already in the
Fusco--Julin good class, the low-gradient part of the level surface is
discarded relative to the intrinsic average
$\bar f=\int f/P$, with $f=|\nabla u|$.  The discarded part is charged
by the Cauchy--Schwarz defect $D_H$ itself.  Consequently the velocity
defect estimate required by the metric first variation holds without any
pointwise lower-gradient hypothesis.  The resulting swept low-gradient
control is only asserted on the isoperimetric-good set $G_I$ used in
\textsc{WeightedMetricTrace}; no unrestricted sweep estimate is claimed.
\end{abstract}

\section{Algebraic truncation on one level}

Let $(S,\Hn)$ be a finite measure space, let $f:S\to(0,\infty)$ be
measurable, and set
\[
   P:=\Hn(S),\qquad
   \alpha:=\int_S \frac{1}{f}\,d\Hn,\qquad
   \beta:=\int_S f\,d\Hn,\qquad
   D:=\alpha\beta-P^2.
\]
Write
\[
   \bar f:=\frac{\beta}{P},\qquad
   \phi:=\frac1f,\qquad
   \bar\phi:=\frac{\alpha}{P}.
\]
By Cauchy--Schwarz, $D\ge0$.

\begin{lemma}[Low-gradient mass is charged by $D$]\label{lem:bad-alpha}
Let
\[
   B:=\{x\in S:f(x)<\bar f/2\},\qquad G:=S\setminus B,
\]
and put $\alpha_B:=\int_B f^{-1}\,d\Hn$.  Then
\begin{equation}\label{eq:alphaB}
   \alpha_B\le \frac{4}{\beta}\,D,
   \qquad
   \Hn(B)\le\frac{\bar f}{2}\,\alpha_B.
\end{equation}
\end{lemma}

\begin{proof}
The weighted variance identity is
\begin{equation}\label{eq:weighted-var}
   \int_S \frac{(f-\bar f)^2}{f}\,d\Hn
   =\beta-2\bar f P+\bar f^2\alpha
   =\frac{\beta}{P^2}D.
\end{equation}
On $B$ one has $\bar f-f>\bar f/2$, hence
\[
   \frac1f\le \frac{4}{\bar f^2}\frac{(f-\bar f)^2}{f}.
\]
Integrating and using $\bar f=\beta/P$ in \eqref{eq:weighted-var} gives
$\alpha_B\le4D/\beta$.  The perimeter estimate follows from
$1\le(\bar f/2)f^{-1}$ on $B$.
\end{proof}

\begin{lemma}[Unit-weight reciprocal variance after truncation]
\label{lem:trimmed-phi}
With $B,G$ as above,
\begin{equation}\label{eq:phi-trimmed}
   \int_S |\phi-\bar\phi|\,d\Hn
   \le \frac{2P}{\beta}\,D^{1/2}
       +\frac{20}{\beta}\,D .
\end{equation}
\end{lemma}

\begin{proof}
Let $\phi_G:=\Hn(G)^{-1}\int_G\phi\,d\Hn$ if $\Hn(G)>0$; if
$\Hn(G)=0$, the estimate below is immediate from Lemma~\ref{lem:bad-alpha}.
On $G$ we have $f\ge\bar f/2$, so
\[
   \bigl(\phi(x)-\phi(y)\bigr)^2
   =\frac{(f(x)-f(y))^2}{f(x)^2f(y)^2}
   \le\frac{4P^2}{\beta^2}
        \frac{(f(x)-f(y))^2}{f(x)f(y)} .
\]
Therefore, by the double-integral variance identity on $G$,
\[
   \left(\int_G|\phi-\phi_G|\,d\Hn\right)^2
   \le \Hn(G)\int_G(\phi-\phi_G)^2\,d\Hn
   \le \frac{4P^2}{\beta^2}D .
\]
It remains to replace the good mean $\phi_G$ by the global mean
$\bar\phi$.  Since $\phi\le2P/\beta$ on $G$ and
$\Hn(B)\le(\bar f/2)\alpha_B$, we have
\[
   \Hn(G)|\phi_G-\bar\phi|
   \le \frac{\Hn(B)}{P}\int_G\phi\,d\Hn
       +\frac{\Hn(G)}{P}\int_B\phi\,d\Hn
   \le 2\alpha_B .
\]
Also,
\[
   \int_B|\phi-\bar\phi|\,d\Hn
   \le \alpha_B+\bar\phi\,\Hn(B)
   \le 3\alpha_B,
\]
because
\[
   \bar\phi\,\Hn(B)
   =\frac{\alpha_G+\alpha_B}{P}\Hn(B)
   \le \frac{\Hn(B)}{P}\frac{2P^2}{\beta}+\alpha_B
   =\frac{\Hn(B)}{\bar f/2}+\alpha_B
   \le 2\alpha_B .
\]
Combining these estimates gives
\[
   \int_S|\phi-\bar\phi|\,d\Hn
   \le \frac{2P}{\beta}D^{1/2}+5\alpha_B,
\]
and Lemma~\ref{lem:bad-alpha} gives \eqref{eq:phi-trimmed}.
\end{proof}

\section{Application to torsion level surfaces}

Let $\Om\subset B_R$ have $|\Om|=\omega_n$, let $u=u_\Om$, and write
$E_\rho=\{u>t(\rho)\}$ with $|E_\rho|=\omega_n\rho^n$.
We use the term \emph{admissible regular radius} for a radius
$\rho\in(0,1)$ satisfying all of the following standard a.e. properties:
$t$ is differentiable at $\rho$; $E_\rho$ has finite perimeter;
$S_\rho:=\partial^*E_\rho$ is a reduced-boundary representative for the
coarea level; the trace of $|\nabla u|$ on $S_\rho$ is finite and
strictly positive $\Hn$-a.e.; and the coarea derivative and flux
identities below hold.  The set of such radii has full Lebesgue measure
in $(0,1)$, after the usual exclusion of exceptional torsion levels.
At an admissible regular radius we set
\[
   S_\rho:=\partial^*E_\rho,\qquad
   f:=|\nabla u|,\qquad
   P:=P(E_\rho),\qquad
   \beta:=\int_{S_\rho}f\,d\Hn=|E_\rho|=\omega_n\rho^n.
\]
Let
\[
   D_H(t(\rho))
   =\left(\int_{S_\rho}\frac1f\,d\Hn\right)
      \left(\int_{S_\rho}f\,d\Hn\right)-P(E_\rho)^2.
\]
The normal velocity is
\[
   V_\rho=\frac{-t'(\rho)}{f},
   \qquad
   \bar V_\rho=\frac{1}{P(E_\rho)}\int_{S_\rho}V_\rho\,d\Hn
              =\frac{P(B_\rho)}{P(E_\rho)}.
\]
The identity
\begin{equation}\label{eq:lambda-alpha}
   (-t'(\rho))\int_{S_\rho}\frac1f\,d\Hn=P(B_\rho)
\end{equation}
is the volume-radius derivative of $|E_\rho|=\omega_n\rho^n$ and is used
only at admissible regular radii.

\begin{proposition}[Velocity defect without Hyp-G]\label{prop:no-hypg}
Fix $\rho_c\in(0,1)$ and $P_1<\infty$.  There is
$C=C(n,\rho_c,P_1)$ such that, for every admissible regular
$\rho\in[\rho_c,1]$ satisfying $P(E_\rho)\le P_1$,
\begin{equation}\label{eq:no-hypg}
   \left(\int_{\partial^*E_\rho}
      |V_\rho-\bar V_\rho|\,d\Hn\right)^2
   \le C(n,\rho_c,P_1)\,D_H(t(\rho)).
\end{equation}
No pointwise lower bound on $|\nabla u|$ is assumed.  The only lower
nondegeneracy input is the volume lower bound
$\beta=|E_\rho|\ge\omega_n\rho_c^n$, and the only upper geometric input
is the perimeter bound $P(E_\rho)\le P_1$.
\end{proposition}

\begin{proof}
Apply Lemma~\ref{lem:trimmed-phi} on $S_\rho$ with
$D=D_H(t(\rho))$.  Since $\rho\ge\rho_c$ and $P(E_\rho)\le P_1$,
\[
   \beta=\omega_n\rho^n\ge\omega_n\rho_c^n=:\beta_0,
   \qquad P/\beta\le P_1/\beta_0 .
\]
Let $\lambda:=-t'(\rho)\ge0$.  Then
\[
   \int_{S_\rho}|V_\rho-\bar V_\rho|\,d\Hn
   =\lambda\int_{S_\rho}|\phi-\bar\phi|\,d\Hn
   \le C_1\lambda D_H^{1/2}+C_2\lambda D_H .
\]
Talenti's comparison gives $\lambda\le1/n$.  For the quadratic term,
use $D_H\le\alpha\beta$ and \eqref{eq:lambda-alpha}:
\[
   \lambda^2D_H
   \le \lambda^2\alpha\beta
   =\lambda\,P(B_\rho)\,\beta
   \le C(n).
\]
Therefore
\[
   \left(C_1\lambda D_H^{1/2}+C_2\lambda D_H\right)^2
   \le C\,D_H,
\]
which proves \eqref{eq:no-hypg}.
\end{proof}

\begin{corollary}[Fusco--Julin good levels provide $P_1$]
\label{cor:FJ-good}
Let
\[
   G_I:=\{\rho: D_I(t(\rho))\le\eta_I\}.
\]
If $\eta_I=\eta_I(n,\rho_c)$ is chosen fixed and small, then
$P(E_\rho)\le P_1(n,\rho_c)$ on $G_I\cap[\rho_c,1]$.  Hence
\eqref{eq:no-hypg} holds for a.e. admissible Fusco--Julin good level in
the collar $[\rho_c,\rho_\delta]$.
\end{corollary}

\begin{proof}
Since
\[
   D_I(t(\rho))=P(E_\rho)^2-P(B_\rho)^2,
\]
the condition $D_I(t(\rho))\le\eta_I$ gives
$P(E_\rho)^2\le P(B_\rho)^2+\eta_I\le n^2\omega_n^2+\eta_I$ for
$\rho\le1$.  This gives a uniform $P_1$.
\end{proof}

\section{Discarded low-gradient sweep on $G_I$}

The same algebra makes precise the idea of discarding tiny regions where
the gradient is anomalously small, but only after restricting to the
isoperimetric-good set.  The estimate below is not a global sweep
estimate over all levels.

\begin{corollary}[The good-level discarded sweep is deficit-controlled]
\label{cor:sweep}
On an admissible good level $\rho\in G_I\cap[\rho_c,\rho_\delta]$ let
\[
   \bar f_\rho:=\frac{|E_\rho|}{P(E_\rho)},\qquad
   B_\rho^{\rm low}:=\{x\in\partial^*E_\rho:
      |\nabla u(x)|<\bar f_\rho/2\}.
\]
Then
\begin{equation}\label{eq:low-velocity}
   \int_{B_\rho^{\rm low}}V_\rho\,d\Hn
   \le C(n,\rho_c)\,(-t'(\rho))\,D_H(t(\rho)).
\end{equation}
Consequently,
\begin{equation}\label{eq:swept-volume}
   \int_{G_I\cap[\rho_c,\rho_\delta]}
      \int_{B_\rho^{\rm low}}V_\rho\,d\Hn\,d\rho
   \le C(n,\rho_c)\,\delta_T(\Om).
\end{equation}
The outer integral is understood after intersecting $G_I$ with the
full-measure set of admissible regular radii.
\end{corollary}

\begin{proof}
Equation \eqref{eq:low-velocity} is Lemma~\ref{lem:bad-alpha} multiplied
by $\lambda=-t'(\rho)$, using
$\beta=|E_\rho|\ge\omega_n\rho_c^n$.  Integrating in $\rho$ gives
\[
   \int_{G_I\cap[\rho_c,\rho_\delta]}
      (-t'(\rho))D_H(t(\rho))\,d\rho
   =\int_{G_I\cap[\rho_c,\rho_\delta]}D_H(t(\rho))\,d\mu(\rho),
\]
which is bounded by $C(n,\rho_c)\delta_T(\Om)$ by the exact level-set
deficit identity on the collar.
\end{proof}

\section{Replacement in the metric trace block}

In \textsc{WeightedMetricTrace}, the metric derivative is estimated on
the isoperimetric-good set $G_I$ and at a.e. admissible regular radius by
\[
   |\dot F_\rho|_\X
   \le C_{n,\rho_c}\int_{\partial^*E_\rho}
      |V_\rho-H_{z_\rho,\rho}|\,d\Hn .
\]
The homothetic trace term is already bounded by the Fusco--Julin
radial trace:
\[
   \left(\int_{\partial^*E_\rho}
      |H_{z_\rho,\rho}-\bar V_\rho|\,d\Hn\right)^2
   \le C\,D_I(t(\rho)).
\]
Proposition~\ref{prop:no-hypg} replaces the old Hyp-G estimate for the
velocity term:
\[
   \left(\int_{\partial^*E_\rho}
      |V_\rho-\bar V_\rho|\,d\Hn\right)^2
   \le C\,D_H(t(\rho)).
\]
This import uses exactly the two hypotheses verified above on $G_I$:
the volume lower bound $|E_\rho|\ge\omega_n\rho_c^n$ and the
perimeter good-level bound $P(E_\rho)\le P_1(n,\rho_c)$.  It does not
use, assume, or prove any pointwise lower bound for $|\nabla u|$.
Therefore, for a.e. $\rho\in G_I\cap[\rho_c,\rho_\delta]$,
\[
   |\dot F_\rho|_\X^2
   \le C(n,R,\rho_c)\bigl(D_H(t(\rho))+D_I(t(\rho))\bigr).
\]
On the complementary isoperimetric-bad set $B_I$, the existing proof uses
the uniform metric Lipschitz bound (A0) and
$\mu(B_I)\le C\delta_T(\Om)$; this part is unchanged.  Thus the
integrated action estimate of \textsc{WeightedMetricTrace} is recovered
without Hyp-G(lower).

\begin{remark}[Collar choice]
The proof only requires a fixed lower radius $\rho_c>0$ to keep
$|E_\rho|$ bounded below and to give an order-one trace window.  Since
the final boundary-layer transfer uses the endpoint
$\rho_\delta=1-\kappa\sqrt{\delta_T}$, one may choose, for instance,
$\rho_c=3/4$ once $\delta_T$ is small.  This makes the repair genuinely
collar-local near the outer boundary while preserving an order-one
window for the Markov plus kinetic transport argument.
\end{remark}

\begin{remark}[Status]
This note does not prove a pointwise lower gradient bound.  It shows that
such a bound is unnecessary for the Plan 2 metric trace: the portion of a
good admissible level where $|\nabla u|$ is too small is self-truncated
and paid for by $D_H$.  Tiny remote components are therefore harmless at
this good-level step, because their $G_I$-restricted low-gradient sweep
contributes to the defect budget rather than to a uniform lower-gradient
constant.
\end{remark}

\end{document}
